\subsection{Log-gamma integrals: the missing even-character contribution}
\label{corr:loggamma}

The abstract and Section~5 of
\path{polylog-polygamma-bridge.tex} suggest that
$\psi^{(-2)}(a)=\int_0^a\log\Gamma(t)\,dt$ at every rational $a$ reduces
to weight-two Clausen values, $\zeta'(-1)$, and logarithms.  The three
displayed examples, at denominators $3,4,6$, are correct.  Their
extrapolation to all denominators omits an even-character contribution.
The same extrapolation occurs in Section~1 of
\path{loggamma-integrals-clausen-bridge__c5013328db83.md}.
Here is an explicit correction at denominator $5$.

\begin{proposition}[An exact denominator-five correction]
\label{corr:loggamma-five}
Let $\chi_5$ be the quadratic character modulo $5$, with values
$(1,-1,-1,1)$ on residues $1,2,3,4$.  Then
\begin{align}
 \int_0^{1/5}\log\Gamma(t)\,dt
 ={}&\frac{2}{25}+\frac{\log(2\pi)}{10}
       +\frac{\operatorname{Cl}_2(2\pi/5)}{4\pi}
       -\frac65\zeta'(-1) \notag\\
    &+\frac1{20}L'(-1,\chi_5)-\frac{29}{1200}\log5.
 \label{corr:loggamma-five-formula}
\end{align}
More generally, for $a>0$,
\begin{equation}\label{corr:loggamma-general}
 \int_0^a\log\Gamma(t)\,dt
 =\zeta'(-1,a)-\zeta'(-1)
       +\frac{a-a^2}{2}+\frac a2\log(2\pi).
\end{equation}
\end{proposition}
\begin{proof}
Differentiate $\partial_a\zeta(s,a)=-s\zeta(s+1,a)$ with respect to
$s$ at $s=-1$.  Lerch's identity and $\zeta(0,a)=\tfrac12-a$ give
\[
 \partial_a\zeta'(-1,a)
 =\zeta'(0,a)-\zeta(0,a)
 =\log\Gamma(a)-\tfrac12\log(2\pi)-\tfrac12+a.
\]
The shift identity for Hurwitz zeta shows that
$\zeta'(-1,a)\to\zeta'(-1)$ as $a\downarrow0$, since the additional
term is $-a\log a$.  Integration proves
\eqref{corr:loggamma-general}.  These are classical Hurwitz-zeta
identities \cite{DLMF25}; the integrated theory is also developed by
Espinosa and Moll \cite{EspinosaMoll2002b}.

Put $Z_r=\zeta'(-1,r/5)$ and $Z=\zeta'(-1)$.  The Fourier reflection
formula at this weight is
\begin{equation}\label{corr:zeta-minus-one-reflection}
 \zeta'(-1,a)-\zeta'(-1,1-a)
       =\frac{\operatorname{Cl}_2(2\pi a)}{2\pi},
       \qquad 0<a<1.
\end{equation}
For completeness, subtract the two Hurwitz Fourier expansions
\[
 \zeta(1-s,a)-\zeta(1-s,1-a)
 =\frac{4\Gamma(s)}{(2\pi)^s}\sin\frac{\pi s}{2}
       \sum_{m\geq1}\frac{\sin(2\pi ma)}{m^s}
\]
and differentiate at $s=2$.  Absolute local convergence of the
series and its logarithmic derivative justifies the differentiation;
the zero of the sine factor leaves precisely
\eqref{corr:zeta-minus-one-reflection}.

Differentiating distribution and the character decomposition at $s=-1$
gives
\begin{align*}
 Z_1+Z_2+Z_3+Z_4&=-\frac45 Z-\frac{\log5}{60},\\
 Z_1-Z_2-Z_3+Z_4
   &=\frac15\bigl(L'(-1,\chi_5)+\log5\,L(-1,\chi_5)\bigr).
\end{align*}
Here $L(-1,\chi_5)=-2/5$, obtained directly from
$\zeta(-1,a)=-B_2(a)/2$.  Combining these equations with
$Z_1-Z_4=\operatorname{Cl}_2(2\pi/5)/(2\pi)$ yields
\[
 Z_1=-\frac15 Z+\frac1{20}L'(-1,\chi_5)
       -\frac{29}{1200}\log5
       +\frac{\operatorname{Cl}_2(2\pi/5)}{4\pi}.
\]
Insert this into \eqref{corr:loggamma-general} at $a=1/5$.
\end{proof}

The correction concerns the scope of the available reduction laws.
It does not prove that $L'(-1,\chi_5)$ is transcendental or independent
of the other constants in \eqref{corr:loggamma-five-formula}.  A precise
replacement for the earlier assertion is: the odd reflection component
is Clausen-explicit, while the even component requires the appropriate
Dirichlet-$L$ derivatives unless further relations have been proved.
The later \path{stieltjes-antiderivative-ladder.tex} already gives the
general Hurwitz-jet description; the earlier bridge discussion should
be brought into agreement with it.

\subsection{Two Herglotz classification statements have exact counterexamples}
\label{corr:herglotz}

In this subsection $F$ denotes the Herglotz--Zagier function, and
$J$ its integral companion:
\[
 F(x)=\sum_{m\geq1}\frac{\psi(mx)-\log(mx)}{m},\qquad
 J(x)=\int_0^1\frac{\log(1+t^x)}{1+t}\,dt,\qquad x>0.
\]
Their classical relation is
\begin{equation}\label{corr:J-F}
 J(x)=F(2x)-2F(x)+F(x/2)+\frac{\pi^2}{12x}.
\end{equation}
See Radchenko and Zagier \cite{RadchenkoZagier2023} for these functions,
the functional equations, and the finite dilogarithm representation.
They are unrelated to the function $F_{n,k}$ in
Section~\ref{zero:section}.

The proposition labelled \texttt{prop:collapse} in
\path{herglotz-arithmetic-study__467d9018ec87.tex} asserts a criterion
involving $\operatorname{rad}(q)\mid30$ for reducing $F(p/q)-F(1)$
to rational-argument dilogarithms, logarithms, and $\pi^2$.
Its necessity fails already at $p/q=1/7$.  Indeed, the report's own
proved formula \texttt{thm:F1q} is
\begin{equation}\label{corr:F-unit-fraction}
 F(1/q)-F(1)
 =-\frac{(q-1)(3q-2)}{24q}\pi^2-\frac12 S(q),
 \qquad
 S(q)=\sum_{r=1}^{q-1}\log^2\left(2\sin\frac{\pi r}{q}\right).
\end{equation}
In particular,
\begin{equation}\label{corr:F-seventh}
 F(1/7)-F(1)
   =-\frac{19}{28}\pi^2
      -\sum_{r=1}^{3}\log^2\left(2\sin\frac{\pi r}{7}\right).
\end{equation}
This contains no dilogarithm at all, although $7\nmid30$.

To check the mechanism independently, specialize the finite
root-of-unity dilogarithm formula to numerator $1$:
\[
 F(1/q)-F(1)=
 \sum_{\substack{\beta^q=1\\\beta\ne1}}
   \left[\operatorname{Li}_2\left(\frac{\beta}{\beta-1}\right)
          -\frac{\pi^2}{6}\right].
\]
If $\beta=e^{2\pi i r/q}$, the transformed argument has real part
$1/2$, modulus $(2\sin(\pi r/q))^{-1}$, and argument
$\pi(r/q-1/2)$.  Euler reflection, paired with conjugation, evaluates
its real dilogarithm component.  Summing the squared arguments and
treating $\beta=-1$ separately when $q$ is even gives
\eqref{corr:F-unit-fraction}.

Nor does restricting to numerators greater than $1$ repair the criterion.
The same laws give
\begin{align}
 F(6/7)-F(1)={}&-\frac8{63}\pi^2+\operatorname{Li}_2(6/7)
       +\frac12\log^2(7/6)\notag\\
       &-\frac12S(7)+\frac12S(6).
 \label{corr:F-six-sevenths}
\end{align}
For a direct derivation, apply the three-term functional equation to
$x=1/6$ and the inversion formula to $x=7/6$:
\begin{align*}
 F(x)-F(x+1)-F\left(\frac{x}{x+1}\right)
   &=-F(1)+\operatorname{Li}_2\left(\frac1{x+1}\right),\\
 F(x)+F(1/x)-2F(1)
   &=\frac12\log^2x-\frac{\pi^2(x-1)^2}{6x}.
\end{align*}
Eliminate $F(7/6)$ and use \eqref{corr:F-unit-fraction} for $q=6,7$.
Thus the appropriate correction is to remove the asserted
classification and retain the proved families and individually
established reductions.  These counterexamples do not determine a
replacement classification for all numerators and denominators.

The proposition labelled \texttt{prop:J} in the same report asserts
that $J(2/q)$ has the specified elementary form only for $q=3,5$.
The allowed elementary inventory in its examples includes $\pi^2$.
With that convention, $q=4$ is another exact counterexample:
\begin{equation}\label{corr:J-half}
 \boxed{\displaystyle J(1/2)=\frac{\pi^2}{48}
                                  +\frac14\log^22.}
\end{equation}
Indeed, \eqref{corr:F-unit-fraction} gives
\[
 F(1/2)-F(1)=-\frac{\pi^2}{12}-\frac12\log^22,
 \qquad
 F(1/4)-F(1)=-\frac{5\pi^2}{16}-\frac34\log^22.
\]
Their substitution in \eqref{corr:J-F} proves \eqref{corr:J-half}.
It also disproves the following sentence in that proposition, which
claims that $J(1/q)$ always retains dilogarithms.  Adding an odd-$q$
restriction would avoid this particular example, but would not prove
the remaining necessity claim.

\subsection{The even-denominator term in the Herglotz derivative formula}
\label{corr:herglotz-derivative}

The same report's \texttt{thm:deriv} defines
\[
 W(q)=\sum_{r=1}^{q-1}
       \cot\frac{2\pi r}{q}\,
       \operatorname{Cl}_2\left(\frac{2\pi r}{q}\right).
\]
For even $q$ the displayed summand at $r=q/2$ is undefined as written:
$\cot\pi$ is singular and $\operatorname{Cl}_2(\pi)=0$.
Define instead
\begin{equation}\label{corr:W-definition}
 W_*(q)=\sum_{\substack{1\leq r<q\\2r\ne q}}
       \cot\frac{2\pi r}{q}\,
       \operatorname{Cl}_2\left(\frac{2\pi r}{q}\right).
\end{equation}
The corrected formula, retaining the report's separate correction, is
\begin{equation}\label{corr:F-derivative}
 \frac2qF'(2/q)-(1+\log2)
    =\frac{q^2-4}{24q}\pi^2+W_*(q)
                                  -\mathbf1_{2\mid q}\log2.
\end{equation}
Equivalently, extend the product continuously at angle $\pi$ and
include it in the sum, while removing the separate $-\log2$ term.
These conventions must not be combined.

The continuous value follows from
$\operatorname{Cl}_2'(\theta)=-\log(2\sin(\theta/2))$ near $\pi$:
\[
 \operatorname{Cl}_2(\pi+u)=-u\log2+O(u^3),\qquad
 \cot(\pi+u)=u^{-1}+O(u),
\]
so the product tends to $-\log2$.  This explains exactly the finite
term required in \eqref{corr:F-derivative}.
At $q=4$ both remaining nonzero-angle cotangents vanish, so
$W_*(4)=0$ and
\begin{equation}\label{corr:F-derivative-half}
 F'(1/2)=2+\frac{\pi^2}{4}.
\end{equation}
Consequently the accompanying statement that this derivative family
never reduces to elementary form is also too broad.

\subsection{Cubic subfields, class fields, and the regulator coefficient}
\label{corr:stark-fields}

The report
\path{herglotz-stark-regulators__6096eccfa2a0.tex}, in its discussion of
$\mathbb Q(\sqrt{257})$ and in the subsequent table, uses the same
symbol $H_0$ both for a cubic field over $\mathbb Q$ and for a cyclic
cubic class-field extension of a quadratic field.  These have different
absolute degrees.  A degree-three field over $\mathbb Q$ cannot
contain a quadratic field, by the tower law.

In the intended $S_3$ situation the notation should distinguish
\[
 [H:K]=3,\quad[K:\mathbb Q]=2,\quad[H:\mathbb Q]=6,
 \qquad [H_0:\mathbb Q]=3.
\]
Here $H/K$ is the cyclic cubic Hilbert or ring class field extension,
and $H_0$ is a non-Galois cubic subfield of its normal closure $H$.
The polynomial displayed for $H_0$ specifies the latter field, not
$H/K$.  This is the standard class-field/cubic-subfield correspondence
\cite{MilneCFT}; it is compatible with the framework of
\cite{RadchenkoZagier2023}.

There is a related omitted factor in the general explanation of the
leading $L$-coefficient.  Let $\chi$ be either nontrivial character of
$\operatorname{Gal}(H/K)\simeq C_3$.  The two conjugate induced
characters are the standard two-dimensional representation of $S_3$.
Artin formalism therefore gives
\begin{equation}\label{corr:artin-cubic}
 L_K(s,\chi)=\frac{\zeta_{H_0}(s)}{\zeta(s)}.
\end{equation}
One can see the representation identity directly: the permutation
representation on the three embeddings of $H_0$ is the direct sum
of the trivial representation and that standard representation.
If $H_0$ is totally real, the analytic class number formula at zero
then gives \cite{MilneCFT}
\[
 \zeta_{H_0}(s)=-\frac{h(H_0)\operatorname{Reg}(H_0)}{2}s^2+O(s^3),
 \qquad \zeta(0)=-\frac12,
\]
and hence
\begin{equation}\label{corr:leading-cubic-L}
 L_K^*(0,\chi):=\lim_{s\to0}s^{-2}L_K(s,\chi)
            =h(H_0)\operatorname{Reg}(H_0).
\end{equation}
Thus the assertion $L_K^*(0,\chi)=\operatorname{Reg}(H_0)$ needs
$h(H_0)=1$.  The printed general explanation omits that hypothesis.
For the five cubics actually displayed, however, it can be supplied
by a short exact argument.  This repairs the stated leading-coefficient
formula for these fields while distinguishing it from the general
identity \eqref{corr:leading-cubic-L}.  Neither this argument nor the
identity proves the report's separate fitted Herglotz evaluations.
The leading-coefficient identity follows from Artin formalism and the
analytic class number formula; it does not need a new appeal to a
conjectural Stark statement.

\begin{proposition}[Class-number certificates for the five cubic fields]
\label{corr:cubic-class-numbers}
Each of the five totally real cubic fields in
Table~\ref{corr:classnumber-table} has class number one.  In the table,
$\alpha$ denotes a root of the polynomial in its row, $d_H$ is the
field discriminant, and $B=\lfloor 2\sqrt{d_H}/9\rfloor$ is the
integer Minkowski bound.  The final column lists a generator of a
principal ideal of norm $p$ for each indicated rational prime $p$;
the field norm of the generator may be $-p$.
\end{proposition}

\begin{table}[htbp]
\centering
\small
\begin{tabular}{c|r|r|l}
 Polynomial for $\alpha$ & $d_H$ & $B$ & $p:\ \text{generator}$\\\hline
 $X^3-X^2-3X+1$ & $148$ & $2$ & $2:\alpha-1$\\
 $X^3-X^2-5X-1$ & $404$ & $4$ &
       $2:\alpha+1;\quad3:\alpha^2-2\alpha-2$\\
 $X^3-X^2-7X-3$ & $788$ & $6$ &
       $2:\alpha+1;\quad3:\alpha;\quad5:\alpha^2-2\alpha-4$\\
 $X^3-X^2-4X+3$ & $257$ & $3$ & $3:\alpha$\\
 $X^3-X^2-11X-7$ & $2708$ & $11$ &
       $2:\alpha+1;\quad3:\alpha+2;\quad7:\alpha$\\
 & & & $5:-\alpha^2+2\alpha+10;\quad11:2\alpha^2-4\alpha-17$
\end{tabular}
\caption{Exact generators used in the class-number-one proof.}
\label{corr:classnumber-table}
\end{table}

\begin{proof}
The rational-root test proves irreducibility of each polynomial $f$.
Their polynomial discriminants are, in the displayed order,
$148,404,788,257,2708$, all positive.  Thus all five fields are
totally real.  At discriminant $257$, squarefreeness immediately gives
$\mathcal O_{H_0}=\mathbb Z[\alpha]$.  In each other row only $2$
could divide the index.  In all four cases
$\overline f=(X+1)^3$ modulo $2$, while
\[
 \left.\frac{f(X)-(X+1)^3}{2}\right|_{X=1}
                  \in\{-5,-7,-9,-13\}.
\]
These values are odd.  Dedekind's index criterion therefore excludes
$2$ from the index.  Equivalently, each shifted polynomial $f(X+1)$
is Eisenstein at $2$, so its root generates the local ring of integers
\cite[Proposition~7.55 and Remark~7.56]{MilneANT}.  This proves
maximality and the asserted field discriminants.  The use of Dedekind factorization below is consequently
valid also at $2$.

Minkowski's theorem gives an integral ideal in every ideal class
with norm at most $2\sqrt{d_H}/9$ for a totally real cubic field
\cite{MilneANT}.  It suffices to prove principalness for every prime
ideal whose norm is at most $B$.  In the four even-discriminant rows,
$2$ is totally ramified; the displayed norm-$2$ element generates its
unique prime ideal.  At discriminant $257$, the polynomial is
irreducible modulo $2$, so the prime above $2$ has norm $8>B$.

For each relevant odd rational prime $p\leq B$, direct factorization
of $f$ modulo $p$ has type $1+2$.  There is therefore exactly one
prime ideal of norm $p$.  The corresponding element in the table
has norm $\pm p$, so this degree-one prime is principal.  Its
degree-two partner is principal as well, since their product is
$(p)$.  This last observation covers in particular the prime of
norm $9$ in the field of discriminant $2708$.  The table accounts
for all rational primes $p\leq B$, hence for all prime ideals within
the Minkowski bound.  Every ideal class is therefore trivial.

All finite calculations in this proof---discriminants, index tests,
modular splitting types, and generator norms---are reproduced with
integer arithmetic by \path{code/verify_cubic_class_numbers.py}.
\end{proof}

\subsection{Gamma relations and the scope of exact certificates}
\label{corr:gamma-certificates}

The intake README already identifies a major attribution correction
in \path{gamma-lattice-and-certificates.tex}.  Its \texttt{law:ko}
and corresponding abstract claim identify completeness of all
multiplicative relations among rational gamma values with a theorem
of Koblitz--Ogus.  The sufficient algebraicity criterion is a theorem;
the assertion that every algebraic gamma monomial follows from the
standard relations is Rohrlich's conjecture.  The distinction is
explicit, for example, in the gamma-value discussion of
Rivoal and Roques \cite{RivoalRoques2014}.

The exact reductions already obtained from reflection and Gauss
multiplication remain valid.  Their row reduction gives completeness
within the stated formal relation space.  Describing the remaining
numerical gamma values as irreducible among all possible algebraic
monomial relations requires the conjecture.  Accordingly,
\texttt{law:ko} should become a conjecture or a conditional statement,
and the same qualification must be propagated to the abstract,
introduction, and basis claims.  A negative integer-relation search
cannot supply the missing converse.

There is a separate issue in the certificate model.  The article
explains that its principal-branch functional-equation instances
are screened by a numerical branch-domain guard.  Exact cancellation
of all special-function atoms verifies the algebraic replay, but
proves the target identity only when the instantiated functional
identities and their branch conditions are valid.  A finite-precision
small-residual check is not a proof of those conditions.  A complete
certificate should therefore include an exact proof of each relevant
branch condition, or explicitly list that condition as an assumption.
For algebraic substitution points, exact real-algebraic sign tests
and a specified argument sector can often provide this missing part.
Single-valued Bloch--Wigner identities offer another route for claims
that can be expressed entirely in that language.

This qualification does not assert an error in any particular
certified value.  The checker and identity stores themselves were
not included in the inspected directory, so their actual soundness
cannot be audited from these reports alone.  It specifies the proof
obligation left by the model described in the article.
